% ===== 原文 pp.14–15 =====
% tr-02a：Zograf 猜想、体积递推的陈述（p.14）；第 2 讲开头与正交边界测地线（p.15）
% 本片无编号陈述、无图、无编号公式。
% p.15 的 "Lecture 2: The McShane-Mirzakhani identity." 是行内小标题 → \para，禁止写 \section。

% [原文 p.14]
Zograf 猜想提示了体积的一种精细结构
\[
V_{g,n}=(4\pi^2)^{2g+n-3}(2g+n-3)!\frac{1}{\sqrt{g\pi}}\left(1+\frac{c_n}{g}+O(1/g^2)\right),
\]
其中 $n$ 固定、$g$ 趋于无穷。

% [Tutor Note: p.14 的 Zograf 猜想给出的是 \(n\) 固定、\(g\to\infty\) 时的渐近：
%  \(V_{g,n}\sim(4\pi^2)^{2g+n-3}(2g+n-3)!/\sqrt{g\pi}\)，注意指数是 \(2g+n-3\) 而阶乘也是 \((2g+n-3)!\)，
%  两者必须同指数，否则量纲对不上（体积多项式的总次数为 \(6g-6+2n\)，而 \(c_n\) 只依赖 \(n\)）。
%  一个可验算的数值锚点：由下文递推的初始值 \(V_{1,1}(L)=\pi^2/12+L^2/48\) 取 \(L=0\) 得
%  \(V_{1,1}(0)=\pi^2/12\approx0.8225\)，而一次穿孔环面模空间 \(\mathcal{M}_{1,1}\) 的 WP 体积是
%  \(\zeta(2)=\pi^2/6\approx1.6449\)，恰为前者的 2 倍；
%  这个因子 2 就是 p.19 所说"环面上 \(\MCG\) 在 \(\cT\) 上的作用不忠实、涉及一个重数"的来源，
%  后文定理 9 的余体积公式（p.21）还要再遇到它。]

作为体积与相交数递推方法的一个应用，Do 导出了体积多项式的\emph{移除边界}（remove a boundary）关系
\[
\frac{\partial V_{g,n+1}}{\partial L_{n+1}}(L,2\pi\ii)=2\pi\ii(2g-2+n)V_{g,n}(L),\qquad[\text{Do08}].
\]
该关系给出紧致情形的体积 $V_g$。

% [Tutor Note: 这条"移除边界"关系的次数记账可以直接自检：\(V_{g,n}\) 作为边界长度多项式的总次数是
%  \(6g-6+2n\)，右端 \(2\pi\ii(2g-2+n)V_{g,n}(L)\) 的 \(L\)-权次正是这个数；左端 \(V_{g,n+1}\) 的总次数为
%  \(6g-4+2n\)，而 \(\frac{\partial}{\partial L_{n+1}}L_{n+1}^{2k}\big|_{L_{n+1}=2\pi\ii}=2k(2\pi\ii)^{2k-1}\)
%  与 \(L\) 无关，所以求导并代入常值后剩下的 \(L\)-次数不超过 \(6g-6+2n\)，两端一致。
%  注意 \(2g-2+n\) 不是总次数的一半（那一半是 \(3g-3+n\)），别混。]

\para{体积递推的陈述} [Mir07a, 第 5 节]。模空间 $\Tg(L_1,\dots,L_n)/\MCG$ 的 WP 体积 $V_g(L_1,\dots,L_n)$ 是边界长度的一个对称函数，具体如下。

\begin{itemize}[leftmargin=1.6em,itemsep=2pt,topsep=2pt]
\item 对 $L_1,L_2,L_3\ge 0$，形式地置
\[
V_{0,3}(L_1,L_2,L_3)=1
\]
以及
\[
V_{1,1}(L_1)=\frac{\pi^2}{12}+\frac{L_1^2}{48}.
\]
\item 对 $L=(L_1,\dots,L_n)$，记 $\what{L}=(L_2,\dots,L_n)$；当 $(g,n)\neq(1,1)$ 或 $(0,3)$ 时，体积满足
\[
\frac{\partial}{\partial L_1}L_1V_g(L)=\mathcal{A}^{con}_g(L)+\mathcal{A}^{dcon}_g(L)+\mathcal{B}_g(L),
\]
其中
\[
\mathcal{A}^{*}_g(L)=\frac{1}{2}\int_0^\infty\!\!\int_0^\infty\what{\mathcal{A}}^{*}_g(x,y,L)\,xy\dd x\,\dd y
\]
并且
\[
\mathcal{B}_g(L)=\int_0^\infty\what{\mathcal{B}}_g(x,L)\,x\dd x.
\]
\end{itemize}

量 $\what{\mathcal{A}}^{con}_g,\what{\mathcal{A}}^{dcon}_g,\what{\mathcal{B}}_g$ 由函数
\[
H(x,y)=\frac{1}{1+e^{\frac{x+y}{2}}}+\frac{1}{1+e^{\frac{x-y}{2}}}
\]
以及子曲面的模空间体积
\[
\what{\mathcal{A}}^{con}_g(x,y,L)=H(x+y,L_1)V_{g-1}(x,y,\what{L})
\]
以及曲面分解之和

% [Tutor Note: 递推式左边写成 \(\partial_{L_1}\big(L_1V_g\big)\) 而不是 \(\partial_{L_1}V_g\)，是为了让 \(L_1\) 的积分权
%  取 \(L_1\dd L_1\)（对应裤子分解里的 \(xy\dd x\dd y\)）在分部积分后自然出现；右端的 \(\mathcal A^{con}_g,\mathcal A^{dcon}_g,\mathcal B_g\) 都只用到总次数更小的体积：\(V_{g-1}\)、\(V_{g_1}V_{g_2}\)，以及少一条边界的 \(V_g(x,L_2,\dots,\widehat L_j,\dots,L_n)\)，
%  所以这是一条显式公式——对 \(L_1\) 积分一次即可解出 \(V_g\)，不需要移项，这正是"递推"的实现方式。
%  初始值 \(V_{0,3}=1\) 与 \(V_{1,1}(L_1)=\pi^2/12+L_1^2/48\) 是两个基准情形——递推对 \((g,n)=(1,1)\) 与 \((0,3)\) 不适用
%  （这两个基准情形的递推退化、无法再往下分解，故直接给出初值）。
%  可验收的数值锚点：\(V_{1,1}(2\pi)=\frac{\pi^2}{12}+\frac{4\pi^2}{48}=\frac{\pi^2}{6}=\zeta(2)\approx1.6449\)，
%  恰为 \(V_{1,1}(0)=\pi^2/12\approx0.8225\) 的两倍；若把 \(1/48\) 抄成 \(1/24\) 或把 \(2\pi\) 代成 \(\pi\)，这个等式立刻不成立。]

% [原文 p.15]
\[
\what{\mathcal{A}}^{dcon}_g(x,y,L)=\sum_{\substack{g_1+g_2=g\\ I_1\sqcup I_2=\{2,\dots,n\}}}H(x+y,L_1)V_{g_1}(x,L_{I_1})V_{g_2}(y,L_{I_2}),
\]
来定义，其中第二个和只考虑双曲结构之对的分解，且无序集合 $I_1,I_2$ 给出一个划分。第三个量 $\what{\mathcal{B}}_g$ 由和
\[
\frac{1}{2}\sum_{j=1}^n\big(H(x,L_1+L_j)+H(x,L_1-L_j)\big)V_g(x,L_2,\dots,\what{L}_j,\dots,L_n)
\]
定义，其中 $L_j$ 已从 $V_g$ 的自变量表中略去。

% [Tutor Note: p.15 这两个求和有三处极易抄错。①\(\sqcup\) 是无交并，\(I_1,I_2\) 无序，
%  故 \((I_1,I_2)\) 与 \((I_2,I_1)\) 只算一次；②"只考虑双曲结构之对的分解"是说两块子曲面的
%  边界数分别为 \(n_i=|I_i|+1\)，双曲结构存在当且仅当 \(2g_i-2+n_i>0\)，
%  这条判据会排除 \(g_i=0\) 且 \(|I_i|\le1\) 的退化项（例如 \((0,3)\) 裤子要求 \(|I_i|=2\)）；
%  ③\(\what{\mathcal{B}}_g\) 的定义式自带 \(1/2\)，且和号里的两个 \(H\) 是 \(H(x,L_1+L_j)\) 与 \(H(x,L_1-L_j)\)，
%  第二个自变量在 \(L_j>L_1\) 时可负，读作 \(H\) 的定义中 \(e^{(x-y)/2}\) 的通配。]

\emph{基本要点}：体积 $V_g(L_1,\dots,L_n)$ 是用少一对裤子构成的那些曲面的体积的一个适当积分。

\para{第 2 讲：McShane--Mirzakhani 恒等式。}

1991 年，Greg McShane 发现了一次穿孔环面上简单测地线长度之和的一条万有恒等式，[McS98]。该恒等式的一个推广充当映射类群作用的“单位分解”（partition of unity）的类比物。这条恒等式使得可以把作用约化为较小的映射类群的作用。对恒等式的考察始于带测地边界的曲面，以及对从边界到自身的弧的研究。

引入 $\cB$：从边界到边界、相对边界同伦的简单曲线之非平凡自由同伦类所成之集。我们通过考察端点落在同一条边界 $\beta$ 上的简单曲线来说明这一途径；对连接不同边界的简单曲线，分析是类似的。每个同伦类含唯一的一条最短测地线，它在端点处与 $\beta$ 正交——称这些测地线为\emph{正交边界测地线}（ortho boundary geodesics）。若把曲面沿其边界加倍，则正交边界测地线加倍为简单闭测地线。

集合 $\cB$ 与嵌入于曲面中、以 $\beta$ 为一条边界的拓扑裤子之集一一对应——称这些裤子为 $\beta$-袖口裤子（$\beta$-cuff pants）。首先注意：正交边界测地线 $\gamma$（$[\gamma]\in\cB$）的两个端点是互异的。$\gamma\cup\beta$ 的一个小邻域/加厚就是相应的拓扑裤子（pair of pants）。几何上，曲线 $\gamma$ 把 $\beta$ 分成若干真子弧；每条子弧与 $\gamma$ 的并是一条简单曲线，它定义一个含唯一测地线的自由同伦类。相应的几何裤子 $\cP$ 以 $\beta$ 与两条确定的测地线为边界。下面我们将看到，一条几何裤子含唯一的一条正交边界测地线。正交边界测地线与 $\beta$ 的并就是“脊”（spines），即嵌入的几何 $\beta$-袖口裤子的线框（wire frames）。

% [Tutor Note: "ortho boundary geodesics"（正交边界测地线）与下页的 "ortho emanating geodesics"
%  （正交出发的测地线）是两个不同的对象，中译都要带"正交"二字以免混淆：
%  前者是单条最短正交弧（每个 \(\cB\) 中类对应一条），后者是从边界正交出发的极大延拓。
%  "把曲面沿边界加倍后正交边界测地线加倍为简单闭测地线"给出了最省事的自检：
%  正交边界测地线必然是"半条"简单闭测地线，它不与曲面内部相交，否则加倍后自交。]

我们现在描述从 $\beta$ 正交出发的测地线的行为如何确定 $\beta$ 的一个 Cantor 子集。该 Cantor 集的测度将为零，而长度恒等式正是其补区间（complementary intervals）的长度之和。下面的描述遵循 Tan--Wong--Zhang 的分析，[TWZ06]。
